\section{Scaling and a limiting terminal pair}\label{sec:scaling}

We now construct a particular ordinary program on the pair of
Theorem~\ref{thm:finite}.  If that program were infinite, we would obtain
strongly \(\Q\)-factorial auxiliary models carrying nef generalized
adjoints and a decreasing sequence of effective boundaries.  Their limit
will be an ordinary terminal pair.  These auxiliary extractions are used
only to compare the original program with a terminating one; the
original program starts on \(X\).

Fix a smooth Kähler form \(h\) on \(X\), with class \(H\), large enough
that \(K_X+\Delta+H\) is nef.  Throughout this section
\(\mathbf H\) denotes the fixed b-nef datum obtained by pulling back
this very form.  On subsequent models its trace class is denoted by
\(H_T\), as in Lemma~\ref{lem:class-trace}.  We do not replace the
fixed datum by arbitrary cohomologous current representatives.

\subsection{Constructing the scaling program}

\begin{proposition}\label{prop:scaling-construction}
Let \((X,\Delta)\) satisfy the hypotheses of Theorem~\ref{thm:finite}.
There is an ordinary negative-ray program, finite or infinite,
\[
 (X,\Delta)=(X_0,\Delta_0)\dashrightarrow
 (X_1,\Delta_1)\dashrightarrow\cdots,
\]
in the same global Weil-divisor \(\Q\)-factorial compact Kähler
category.  It stops when the ordinary adjoint is nef or a negative ray
has a Mori fibre contraction.  Put
\[
 D_i=K_{X_i}+\Delta_i,\qquad
 L_i(t)=D_i+tH_i,\qquad H_i=H_{X_i},\qquad \lambda_{-1}=1.
\]
At every stage before stopping there is a parameter
\(0<\lambda_i\leq\lambda_{i-1}\) and a contracted extremal ray
\(R_i\) such that
\[
 D_i\cdot R_i<0,\qquad L_i(\lambda_i)\cdot R_i=0,
 \qquad L_i(\lambda_i)\text{ is nef}.
\]
The generalized pairs \((X_i,\Delta_i+t\mathbf H)\) are gklt for
\(0\leq t\leq\lambda_{i-1}\).  A birational step is crepant for the
data at \(t=\lambda_i\), and discrepancies do not decrease across
that step for any \(0\leq t\leq\lambda_i\).
\end{proposition}

\begin{proof}
On \(X_0\), the nef datum descends, so its addition changes no
discrepancies.  Thus all the initial generalized pairs are gklt.
Suppose the assertions have been established through \(X_i\).
If \(D_i\) is nef, stop.  Otherwise define
\[
 \lambda_i=\min\{t\in[0,\lambda_{i-1}]:L_i(t)\text{ is nef}\}.
\]
The set is nonempty by induction and is closed and convex.  Its minimum
is positive because \(D_i\) is not nef.

We first find a ray at this threshold.  Choose
\(t_k\uparrow\lambda_i\) with \(\lambda_i/2<t_k<\lambda_i\).
The generalized cone theorem gives an extremal ray negative for
\(L_i(t_k)\).  Since \(L_i(\lambda_i)\) is nef, that ray is also
negative for \(L_i(\lambda_i/2)\).  The trace \(H_i\) is big by
Corollary~\ref{cor:big-trace}.  Hence
\(\Delta_i+(\lambda_i/2)H_i\) is big, and the cone theorem gives
only finitely many negative rays for this latter generalized pair
\cite[Theorem~1.3]{HX2026}.  One ray \(R_i\) therefore occurs for
infinitely many \(k\).  On this fixed ray, continuity gives
\(L_i(\lambda_i)\cdot R_i=0\).  Moreover
\(H_i\cdot R_i>0\), so \(D_i\cdot R_i<0\).

Apply Proposition~\ref{prop:ordinary-step} to this ordinary negative
ray, and denote its contraction by \(f_i:X_i\to Z_i\).  A fibre
contraction gives the stopping alternative.  Otherwise make the
ordinary divisorial step or flip.  Lemma~\ref{lem:class-transport}
transports the trace \(H_i\) to \(H_{i+1}\) while preserving its
relation to the fixed b-datum.  By Lemma~\ref{lem:pullback-image}, the
class \(L_i(\lambda_i)\) descends to \(Z_i\).  Its descended class
is nef, and its pullback on the new model is
\(L_{i+1}(\lambda_i)\).  Thus the two threshold classes have equal
pullbacks on a common resolution.  Their difference is represented by
the actual exceptional discrepancy divisor, which vanishes by
negativity as in Lemma~\ref{lem:small-comparison}.  The step is
consequently generalized crepant at the threshold.

For completeness, this also proves the induction on singularities.
On a common resolution, the log discrepancy at each prime \(E\) is
affine in \(t\).  Its difference across the step is zero at
\(t=\lambda_i\) and is the ordinary discrepancy increase at zero.
Consequently
\begin{align}
 &a(E,X_{i+1},\Delta_{i+1}+t\mathbf H)
       -a(E,X_i,\Delta_i+t\mathbf H)\notag\\
 &\qquad=\left(1-\frac{t}{\lambda_i}\right)
     \bigl(a(E,X_{i+1},\Delta_{i+1})-a(E,X_i,\Delta_i)\bigr)
     \geq0\label{eq:scaling-interpolation}
\end{align}
for \(0\leq t\leq\lambda_i\).  The generalized pairs on this
interval remain gklt, and the new upper-threshold class is nef.  The
ordinary steps preserve all the stated properties of the original
category by Proposition~\ref{prop:ordinary-step}.
\end{proof}

\subsection{An infinite program has vanishing thresholds}

We record first why an infinite program has a tail consisting only of
flips.  Use the cycle-space notation of Lemma~\ref{lem:cycle-loss}:
\(c_3(T)\) is the dimension of the span of compact analytic
three-dimensional cycle classes in \(H_6(T,\R)\).  It is finite.
A small modification of fourfolds is an isomorphism away from sets of
dimension at most two, so that lemma, applied in both directions,
preserves \(c_3\).  A divisorial contraction strictly lowers it: an
exceptional divisor has a nonzero cycle class, detected by integration
of the cube of a Kähler form, and maps to a set of dimension at most
two.  Therefore only finitely many divisorial steps occur.

\begin{lemma}\label{lem:scaling-zero}
If the program of Proposition~\ref{prop:scaling-construction} is
infinite, then \(\lambda_i\to0\).
\end{lemma}

\begin{proof}
Fix \(0<\epsilon<1\) and a resolution \(\nu:U\to X\).  Consider the
compact segment of generalized data
\[
 \mathcal P_\epsilon
 =\{(\Delta,t\nu^*H):\epsilon\leq t\leq1\}
 \subset \operatorname{Div}_{\R}(X)\times\BC(U).
\]
The upstairs nef data here are the pullbacks of the fixed smooth form.
Every pair in this segment is gklt, its upstairs class is nef, and its
boundary-plus-trace is modified big, since \(tH\) is Kähler on \(X\).
This real segment is allowed in the finiteness theorem for weak log
canonical models \cite[Theorem~5.15]{HX2026}; neither rationality of
\(H\) nor strong \(\Q\)-factoriality of \(X\) is required there.

Let \(\phi_i:X\dashrightarrow X_i\) be the accumulated map, and
suppose \(\lambda_i\geq\epsilon\).  It is a weak log canonical
model of the original generalized pair at parameter \(\lambda_i\).
Indeed, it extracts no divisors, its target adjoint is nef, and
discrepancies do not decrease: apply
\eqref{eq:scaling-interpolation} to every preceding step at the fixed
parameter \(\lambda_i\), which is no larger than any preceding
threshold.  On a common resolution the original adjoint is the
pullback of the target nef class plus an effective exceptional
divisor.  Pseudo-effectivity descends by
Lemma~\ref{lem:positive-descent}, so the original data belong to the
pseudo-effective portion of \(\mathcal P_\epsilon\).

The finiteness conclusion concerns the maps \(\phi_i\), up to
isomorphisms of their targets commuting with the maps from \(X\).
These marked maps are pairwise distinct in the flip tail.  In fact,
given two stages, an exceptional prime has strictly increasing
ordinary discrepancy at the first intervening flip by
Lemma~\ref{lem:negativity}, and later steps
cannot decrease it.  A target isomorphism compatible with the marking
would identify the transformed boundary and all its discrepancies,
a contradiction.  This is also the distinction used in the proof of
\cite[Theorem~5.16]{HX2026}.

There are therefore only finitely many \(i\) with
\(\lambda_i\geq\epsilon\).  Since \(\epsilon>0\) was arbitrary,
the thresholds tend to zero.
\end{proof}

An infinite program would thus provide nef generalized adjoints with
arbitrarily small nef parts.  We next put their nonterminal places on
auxiliary models, so that the boundary can be allowed to converge on
one fixed space.

\subsection{Stabilizing the extractions}

\begin{proposition}\label{prop:terminal-limit}
Suppose the program of Proposition~\ref{prop:scaling-construction} is
infinite.  After discarding finitely many stages, there are proper
bimeromorphic morphisms
\[
 \mu_i:Y_i\longrightarrow X_i
\]
with \(Y_i\) strongly \(\Q\)-factorial and compact Kähler, and
effective real boundaries \(\Theta_i\), with the following properties.
\begin{enumerate}[label=\textup{(\roman*)}]
\item The natural maps between the \(Y_i\) are isomorphisms in
codimension one, and
\begin{equation}\label{eq:extracted-nef-adjoint}
 Q_{i,Y_i}:=K_{Y_i}+\Theta_i+\lambda_iH_{Y_i}
       =\mu_i^*L_i(\lambda_i)
\end{equation}
is nef.  Each ordinary pair \((Y_i,\Theta_i)\) is klt and terminal.
\item On a fixed first model \(Y=Y_{i_0}\), write
\(\Theta_{i,Y}\) for the strict transform of \(\Theta_i\).
These boundaries have a fixed finite containing support and decrease
coefficientwise to an effective real boundary \(\Theta_Y\).
The ordinary pair \((Y,\Theta_Y)\) is klt and terminal.
\item The strict transform of \(\Theta_Y\) on each \(Y_i\)
pushes forward to \(\Delta_i\) under \(\mu_i\).
\end{enumerate}
Here terminal means that every exceptional prime has log discrepancy
strictly greater than one.
\end{proposition}

\begin{proof}
Discard the finitely many divisorial steps.  All remaining maps
\(X_i\dashrightarrow X_{i+1}\) are small.  For each such \(i\), put
\[
 \mathcal E_i=\{E\text{ exceptional over }X_i:
      a(E,X_i,\Delta_i+\lambda_i\mathbf H)\leq1\}.
\]
The fixed nef datum has a carrier projective over \(X_i\): take a
common resolution of the finitely many preceding projective ordinary
step diagrams.  For a flip its graph is projective over either side,
so this construction and a further projective log resolution retain
projectivity over \(X_i\).  On this resolution the datum is the
pullback of \(h\).  Lemma~\ref{lem:low-places} therefore makes
\(\mathcal E_i\) finite.  It allows real coefficients and counts
places of log discrepancy exactly one as well as those below one.

Smallness identifies the primes exceptional over successive
\(X_i\).  Their log discrepancies at the current threshold agree by
crepancy.  On the new model, decreasing the coefficient of the nef
part increases discrepancies: in the trace formula of
Lemma~\ref{lem:class-trace}, the correction divisor is effective.
Thus, for every such prime,
\begin{equation}\label{eq:decreasing-extraction-sets}
 \begin{split}
 a(E,X_{i+1},\Delta_{i+1}+\lambda_{i+1}\mathbf H)
 &\geq a(E,X_{i+1},\Delta_{i+1}+\lambda_i\mathbf H)\\
 &=a(E,X_i,\Delta_i+\lambda_i\mathbf H).
 \end{split}
\end{equation}
It follows that \(\mathcal E_{i+1}\subseteq\mathcal E_i\).
After another truncation they equal one finite set \(\mathcal E\).
Only this set stabilizes; no discreteness of the real discrepancy
values is asserted.

Apply the exact extraction theorem
\cite[Corollary~2.29]{HX2026} to the generalized pair at
\(\lambda_i\) and the prescribed set \(\mathcal E\).
It gives a strongly \(\Q\)-factorial compact Kähler model
\(\mu_i:Y_i\to X_i\) extracting precisely those primes.
Its compact Kähler output is supplied by the relative MMP in its proof.
The input of this theorem need not be
strongly \(\Q\)-factorial, and its nef datum and crepant boundary
may be real.

Let \(\Theta_i\) be the crepant generalized boundary on \(Y_i\).
On the transform of a divisor of \(X_i\), its coefficient is the
coefficient in \(\Delta_i\); on an extracted prime \(E\), it is
\[
 1-a(E,X_i,\Delta_i+\lambda_i\mathbf H)\in[0,1).
\]
Thus \(\Theta_i\) is effective.  The generalized crepant relation
gives \eqref{eq:extracted-nef-adjoint}.  More explicitly, subtract
the ordinary class \(K_{Y_i}+\Theta_i\) from the pulled-back
generalized adjoint and divide by \(\lambda_i>0\).
Lemma~\ref{lem:class-trace} identifies the result as the trace
\(H_{Y_i}\) of our fixed b-datum, with an effective exceptional
correction on a common resolution.  Equation
\eqref{eq:extracted-nef-adjoint} is nef because its right-hand side
is a pullback of a nef class.

Every prime exceptional over \(Y_i\) is exceptional over \(X_i\)
and is absent from \(\mathcal E\), so its generalized log
discrepancy is greater than one.  Dropping the nef part weakly
increases this discrepancy.  Hence \((Y_i,\Theta_i)\) is ordinary
terminal.  Together with the boundary coefficients below one, this
also proves that it is klt.

The flip tail identifies all prime divisors on the \(X_i\), and
each \(Y_i\) extracts exactly the same additional primes.
Consequently the \(Y_i\) are naturally isomorphic in codimension
one.  Fix the first one, \(Y=Y_{i_0}\).  On its primes inherited
from \(X_{i_0}\), the coefficients of \(\Theta_{i,Y}\) are
constant.  On the extracted primes, they decrease by
\eqref{eq:decreasing-extraction-sets}.  They are nonnegative and
have a fixed finite containing support, so their coefficientwise
limit \(\Theta_Y\) exists and
\[
 0\leq\Theta_Y\leq\Theta_{i,Y}\leq\Theta_{i_0,Y}.
\]

Terminality of the limit follows by comparison with the first
boundary on this fixed model.  Since \(Y\) is \(\Q\)-factorial,
the effective real divisor \(\Theta_{i_0,Y}-\Theta_Y\) has
nonnegative order at every divisorial valuation.  Therefore
\[
 a(E,Y,\Theta_Y)\geq a(E,Y,\Theta_{i_0,Y})>1
\]
for every exceptional prime \(E\).  All boundary coefficients
remain below one, so the limit is klt as well.  This argument uses
decrease of the boundary, rather than any assertion that terminality
is closed under limits.  Finally, only coefficients of
\(\mu_i\)-exceptional primes vary.  The transformed limiting
boundary thus pushes forward to \(\Delta_i\), proving (iii).
\end{proof}

The infinite-program assumption has now produced an ordinary terminal
pair on a fixed strongly \(\Q\)-factorial model, together with the
nef adjoints \eqref{eq:extracted-nef-adjoint}.  In the next section we
apply terminal termination to that limiting pair and compare a small
perturbation with one of these nef adjoints.
